\section{Simulation Results}
\label{sec:simulations}

\subsection{Simulation Settings}
\label{sec:simulation-settings}

We consider the satellite uplink in Section~\ref{sec:system} with the
parameters in Table~\ref{tab:simulation-parameters}. Terminal locations
are sampled uniformly over the beam footprint. The propagation delays,
Doppler shifts, and array responses are determined by the same satellite
and terminal geometry. Large-scale propagation effects are compensated.
The channel follows the NTN-TDL-D power delay profile \cite{shen2025}, with the delay
spread drawn uniformly from $[30,50]$~ns. The line-of-sight component has
a uniformly distributed phase and fixed amplitude, while the diffuse
path gains are proper complex Gaussian. Channels remain fixed over each
access block. Unless otherwise stated, the access configuration uses
$K=40$, $Q=16$, and $p=0.1$.

\begin{table}[!t]
\caption{Simulation Parameters}
\label{tab:simulation-parameters}
\centering
\small
\begin{tabular}{@{}ll@{}}
\toprule
Parameter & Value \\
\midrule
Carrier frequency $f_c$ & $2$~GHz \\
Satellite altitude & $600$~km \\
Beam radius & $128$~km \\
Minimum elevation angle & $30^\circ$ \\
UPA dimensions $M_y\times M_z$ & $4\times4$ \\
Array element spacing $d_y,d_z$ & $\lambda/2$ \\
AFDM block length $N$ & $16$ \\
Subcarrier spacing $\Delta f$ & $30$~kHz \\
Chirp parameters $(c_1,c_2)$ & $(7/32,1/320)$ \\
CPP length $N_{\rm CPP}$ & $336$ \\
Residual-delay support & $[0,699]~\mu$s \\
Doppler support & $[-41,41]$~kHz \\
Power delay profile & NTN-TDL-D \\
Delay spread & $[30,50]$~ns \\
Rician factor & $11.707$~dB \\
Default $(K,Q,p)$ & $(40,16,0.1)$ \\
\bottomrule
\end{tabular}
\end{table}

\subsubsection{Channel Samples and Training}
The first-order score network is trained with either
$N_{\rm train}=5{,}000$ or $50{,}000$ clean channel tensors of size
$16\times16\times16$. Four candidate channels are sampled from each
physical scene before activity masking. The smaller training set is a
fixed subset of complete scenes from the larger set. Both models use
the same architecture and a separate validation set of $5{,}000$
channels; training, validation, and test scenes are disjoint.
All channels are expressed in the same fixed training coordinates,
obtained by dividing the physical channel by the fixed reference
scale $0.2647966453$.

We train the first-order network from random initialization using
\eqref{eq:first-loss}. The noise variance is sampled uniformly on a
logarithmic scale over $[10^{-3},10^2]$. Both models use Adam with a
batch size of $64$, an initial learning rate of $10^{-4}$, and an
exponential moving average with decay $0.999$. Training lasts $200$
epochs, and the checkpoint is selected by the minimum fixed-noise
validation loss. The $5{,}000$- and $50{,}000$-sample runs perform
$15{,}800$ and $156{,}400$ optimizer updates, respectively. Thus, the
comparison uses a common epoch budget and different training costs.

\subsubsection{Evaluation Metric}
For channel estimation, we report the normalized mean-squared error
(NMSE), computed by pooling error and channel energies over the test set:
\begin{equation}
 \mathrm{NMSE}_{\rm dB}
 =10\log_{10}\frac{\sum_{i\in\mathcal T}
 \norm{\widehat H_i-H_i}_F^2}
 {\sum_{i\in\mathcal T}\norm{H_i}_F^2},
 \label{eq:simulation-nmse}
\end{equation}
where $\mathcal T$ indexes the evaluated channel samples.
The denoising experiment uses $500$ independent physical scenes with
four channels per scene, giving $2{,}000$ channels at each noise level.
The channels within a scene share its geometry. The curves report
point estimates from this test set.

\subsection{Channel Denoising Performance}
\label{sec:denoising-performance}

We first evaluate the conditional-mean estimator
\eqref{eq:learned-mean} under the direct AWGN model
$R=H+W$, with $\vect(W)\sim\CN(\mathbf0,v\I_D)$ and known $v$.
The nominal denoising SNR is $-10\log_{10}v$ in the fixed training
coordinates, using unit reference channel power. We evaluate
$\{-20,-10,0,10,20,30\}$~dB using the same scene definitions,
candidate channels, and standardized noise draws across methods.
This experiment evaluates the first-order channel mean independently
of activity detection and the receiver iterations.

The reference estimators are least squares (LS), which returns $R$;
the Gaussian-prior estimate $P_{\rm tr}R/(P_{\rm tr}+v)$, where
$P_{\rm tr}$ is the mean element power of the normalized training
channels; and a genie-aided conditional MMSE reference
\cite{cai2026}. The genie is given the true path geometry, path powers,
and line-of-sight complex gain. Its curve uses the analytic conditional
MSE for the remaining Gaussian path gains, normalized by the same
empirical channel energy.
We also evaluate simultaneous orthogonal matching pursuit (SOMP)
\cite{TroppGilbertStrauss2006SOMP}, GAMP-SBL
\cite{ZhangFanHaoQuan2023GAMPSBL}, GNN-ASPL \cite{ding2026}, and
off-grid SBL \cite{WeiEtAl2022OffgridSBL} by applying their
channel-estimation kernels to an identity
observation operator. Their receiver configurations and the GNN weights
are retained, and all methods are supplied with the AWGN variance;
GAMP-SBL retains its internal noise update. These curves therefore
characterize transfer of the existing implementations to direct denoising.

\begin{figure}[!t]
\centering
\includegraphics[width=\columnwidth]{denoising.pdf}
\caption{Independent channel denoising with known AWGN variance:
$500$ physical scenes and four channels per scene. The proposed curves
use $N_{\rm train}=5{,}000$ and $50{,}000$; the four sparse-recovery
baselines retain their receiver configurations. The genie curve uses
true geometry, path powers, and line-of-sight gain.}
\label{fig:denoising-performance}
\end{figure}

Figure~\ref{fig:denoising-performance} shows that the model trained on
$50{,}000$ channels attains lower NMSE than the model trained on
$5{,}000$ channels at all six SNR values, with reductions ranging from
$2.39$ to $7.35$~dB. At $0$~dB, the two NMSE values are $-17.92$
and $-10.57$~dB, respectively. At $30$~dB, the larger model achieves
$-40.22$~dB, compared with $-29.93$~dB for LS. The low-SNR case remains
more challenging: at $-20$~dB, off-grid SBL achieves $-1.81$~dB,
whereas the larger proposed model gives $-0.27$~dB.

The four transferred baselines exhibit an error floor near $-11$~dB.
For SOMP, GAMP-SBL, and GNN-ASPL, the integer delay--Doppler dictionary
limits the channel representation. For off-grid SBL, the retained
implementation estimates coefficients with a first-order atom
approximation and reconstructs the channel with exact atoms, introducing
an additional model discrepancy. Accordingly, the high-SNR gaps in
Fig.~\ref{fig:denoising-performance} include these implementation and
representation effects; they do not establish the performance of
fully adapted versions of the four methods. The proposed estimator
also remains above the genie reference throughout the evaluated range.
